\chapter{一些需要积累的技巧}

\section{行列式求导}

\section{使用矩阵递推数列}

\section{常见恒等式}

\section{不常见的换元技巧}


\chapter{重积分}

\section{很难画出的积分区域}

\begin{example}[紫皮书6.2习题]
已知$f$在$\z[0,1\y]$上连续,证明：
\[ \int_{0}^{1}\mathrm{d} x\int_{0}^{x}\mathrm{d} y\int_{0}^{y}f(x)f(y)f(z)\mathrm{d}z = \frac{1}{3!}\z(\int_{0}^{1}f(x)\mathrm{d}x\y)^3\]
\end{example}

\begin{example}[\;17届CMC]
设 \(\Omega\) 是由平面
$\Pi_1 : x = 0, \Pi_2 : x = y, \Pi_3 : y = z, \Pi_4 : z = 1, \Pi_5 : z = 2$

围成的空间区域，求三重积分$\displaystyle \iiint_{\Omega} \frac{\mathrm{d} x \, \mathrm{d} y\, \mathrm{d} z}{\sqrt{y^2 + z^2}}.$

\end{example}



\chapter{曲面积分}

\section{算子法}

\section{凑微分}

\section{利用物理结论}




\chapter{级数}

\section{倒序求和}

\section{裂项技巧}

\section{二重求和：交换求和次序}